%% notes-tex/034-isobutanol-wetlab/034-isobutanol-wetlab.tex
%% [[experiments.034-isobutanol-wetlab.strain-panel]]
%% https://github.com/Mjvolk3/torchcell/tree/main/notes-tex/034-isobutanol-wetlab/034-isobutanol-wetlab.tex
%%
%% The bench side of isobutanol. Eleven strains from the Avalos lab are in the
%% freezer, and the first question is which of them a first round of RNA
%% sequencing should carry. That question turns on what each strain is, which
%% the tube labels do not say, so the document starts by mapping every tube onto
%% the paper that built it.
%%
%% Build:  make            -> 034-isobutanol-wetlab.pdf       (draft view)
%%         make clean-view -> 034-isobutanol-wetlab-clean.pdf (share view)
%%         make check      -> formatting / provenance / citation gate
%%         make tables     -> regenerate tables/ from the experiment script

\documentclass[10pt,a4paper]{article}
\usepackage[draft]{tcdoc}

\title{\vspace{-1.5em}Isobutanol, Bench Side\\
  {\large Eleven Avalos-lab strains, and which of them a first RNA-seq round carries}}
\author{Michael Volk}
\date{\today}

\begin{document}
\maketitle
\thispagestyle{empty}

\begin{abstract}
\noindent
Eleven strains built for a branched-chain amino acid biosensor are on hand,
labeled by tube number and strain number and by nothing else. Eight of them are
named in the paper that built them, and those eight turn out to be screening
hosts and single-allele controls rather than the strains that produced the
paper's headline titers: three of those headline strains are a single
transformation past a tube already in the freezer. Three more tubes are named
nowhere in the paper's running text at all. For sequencing, the panel spans
roughly an order of magnitude in isobutanol titer, that span is not a clean
ladder because the carbon source changes between arms, and none of these strains
shares a genetic background with the knockout-collection isobutanol work they
get grouped with. A five-strain first round follows from that, anchored on the
screening hosts rather than on a wild type the panel does not contain.
\end{abstract}

\vspace{0.5em}
{\small\color{tcgray}
\textbf{How to read this.} \Cref{sec:panel} is the mapping from tube to strain,
and it is the part that has to be right before anything else is decided.
\Cref{sec:yko} is the relationship to the dissertation work on the yeast
knockout collection, which shares a biosensor cassette with these strains but
not a genetic background. \Cref{sec:rnaseq} is the first sequencing round that
falls out of the two.\par\smallskip

Source paper: Zhang \textit{et al.}, \textit{Nature Communications} \textbf{13},
270 (2022), \url{https://doi.org/10.1038/s41467-021-27852-x}, mirrored as
\file{zhangBiosensorBranchedchainAmino2022}.\par
Generating script: \file{experiments/034-isobutanol-wetlab/scripts/strain\_panel.py}\par}

\vspace{1em}
\tccontents

%% notes-tex/034-isobutanol-wetlab/sections/1-panel.tex
%% SOURCE: experiments/034-isobutanol-wetlab/scripts/strain_panel.py
%%         (tables/t1-strain-panel.tex and tables/t2-one-step-away.tex are emitted
%%         by it; results/strain_panel.csv carries the same records with their quotes)
%% Every genotype and titer traces to the mirrored paper,
%% $DATA_ROOT/torchcell-library/zhangBiosensorBranchedchainAmino2022/.

\section{What Is Actually in the Freezer\secstatus{todo}}
\label{sec:panel}

All eleven tubes come from one paper: Zhang \textit{et al.}, \textit{Nature
Communications} \textbf{13}, 270 (2022). It builds a transcriptional biosensor
for branched-chain higher alcohol (\term{BCHA}) metabolism and then uses it to
run three screens. The sensor is yeast-enhanced green fluorescent protein
(\term{yEGFP}) under the \gene{LEU1} promoter, which Leu3p activates when
$\alpha$-isopropylmalate ($\alpha$-IPM) accumulates, and $\alpha$-IPM tracks
flux through the branched-chain pathway.

\notesec{Two sensor configurations, and they are not interchangeable} The paper
uses the sensor in two forms, and which form a strain carries determines which
product its fluorescence reports.

\begin{quote}\small
the constructs containing yEGFP-PEST/LEU4$_{1-410}$/\gene{leu2}$\Delta$
(henceforth called the isobutanol configuration) and
yEGFP/Leu4$\Delta$S547/\gene{LEU2}/\gene{leu4}$\Delta$/\gene{leu9}$\Delta$
(henceforth called the isopentanol configuration)
\end{quote}

\noindent
The isobutanol configuration uses a destabilized yEGFP (the PEST degron) and a
truncated Leu4p, \term{LEU4$_{1-410}$}, that has lost its leucine regulatory
domain. The isopentanol configuration uses stable yEGFP and a full-length Leu4p
carrying a single serine deletion at position 547,
\term{Leu4$\Delta$S547}. Five of our tubes sit in the first, three in the
second.

\notesec{The background is CEN.PK2-1C, not S288C} This matters more than
anything else in the table, and \cref{sec:yko} is about why.

\begin{quote}\small
\org{S.\ cerevisiae} strain CEN.PK2-1C (MATa \gene{ura3}--52 \gene{trp1}-289
\gene{leu2}-3,112 \gene{his3}-1 MAL2-8c SUC2) and its derivatives were used in
this study
\end{quote}

%% SOURCE: experiments/034-isobutanol-wetlab/scripts/strain_panel.py
%% Mirror: $DATA_ROOT/torchcell-library/zhangBiosensorBranchedchainAmino2022/paper.pdf
%% sha256 551ba08ec3bf584ef49b72186050d4bb3230df59edfa664a8a6af3460a2a8bec
\begin{table}[htbp]
\centering
\caption{\textbf{The eleven strains on hand, and what each one is in the paper that built them.} Color groups are as the tubes are labeled. Titers are not comparable down the column: the carbon source differs by row and two rows are a different product. A dagger marks a value we back-computed from a reported fold change rather than one the paper prints.}
\label{tab:panel}
\footnotesize
\begin{tabular}{@{}llll >{\raggedright\arraybackslash}p{52mm} l@{}}
\toprule
ID & Tube & Group & Config. & What it is & Titer (mg/L) \\
\midrule
C0043 & Yzy311 & green & isobutanol & Strain A of Fig. 3c: the \gene{ILV6} screening host carrying wild-type \gene{ILV6} on a CEN plasmid, with no isobutanol pathway overexpressed. & -- \\
C0044 & YZy313 & green & isobutanol & Strain C of Fig. 3c: wild-type \gene{ILV6} on a CEN plasmid in YZy363, which carries the mitochondrial isobutanol pathway $\delta$-integrated. The denominator of the 1.6-fold improvement. & -- \\
C0045 & YZy314 & green & isobutanol & Strain D of Fig. 3c: the ILV6V110E hit in the same pathway-overexpressing host, 1.6-fold over C0044. One point mutation separates the two. & -- \\
C0046 & YZy148 & blue & isopentanol & Screening host for the \gene{LEU4} mutagenesis libraries. Carries the isopentanol configuration with the \gene{LEU4}$\Delta$S547 allele removed, so a plasmid-borne \gene{LEU4} is the only source of the enzyme. & -- \\
C0047 & YZy148+LEU4WT & blue & isopentanol & Wild-type \gene{LEU4} control for the isopentanol screen: YZy148 retransformed with a plasmid carrying unmutagenized \gene{LEU4}. & 201$^{\dagger}$ \\
C0048 & YZy148+LEU4deltaS547 & blue & isopentanol & The published leucine-insensitive allele, a single Ser547 deletion, carried back into the screening host. This is the benchmark the evolved \gene{LEU4} variants were scored against. & 842 $\pm$ 23 \\
C0049 & YZy452 & orange & isobutanol & Best cytosolic-pathway isolate from three rounds of FACS on $\delta$-site integrants, and the host the \gene{Ll\_ilvD} plasmid library was then screened in. & 310 $\pm$ 15 \\
C0050 & YZy454 & orange & isobutanol & Wild-type \gene{Ll\_ilvD} control: YZy452 carrying a CEN plasmid with the unmutated dehydratase. The denominator of the 3.1-fold improvement. & 73$^{\dagger}$ \\
C0051 & YZy469 & orange & isobutanol & The \gene{Ll\_ilvD} hit, I433V, on a CEN plasmid. Same host and same carbon source as YZy454, so the pair is a clean one-allele contrast. & 227 $\pm$ 13 \\
C0052 & YZy91 & gray & isobutanol & Screening host for the \gene{ILV6} mutagenesis libraries. \gene{BAT1} and \gene{BAT2} are deleted so valine and $\alpha$-KIV cannot interconvert, and \gene{ILV6} is deleted so the only Ilv6p present is the plasmid-borne one under test. The tube itself carries no ILV6; the screen's best hit, this host plus ILV6V110E, reached 378 $\pm$ 10 mg/L. & -- \\
C0053 & YZy502 & gray & isobutanol & Optogenetic host: triple \gene{pdc}$\Delta$ with light-inducible \gene{PDC1} under OptoEXP and the dark-inducible cytosolic pathway under OptoINVRT7, $\delta$-integrated. One plasmid and two FACS rounds short of the paper's best isobutanol strain. & -- \\
\bottomrule
\end{tabular}
\end{table}


\notesec{The tubes are hosts and controls, not the headline strains} Read down
the titer column and the panel looks weaker than the paper does. That is not an
accident of which tubes were shipped. Eight of the eleven are either the host a
mutagenesis library was screened in or the wild-type-allele control that a hit
was scored against. The strains the paper leads with are the hits, and three of
them are one transformation away.

%% SOURCE: experiments/034-isobutanol-wetlab/scripts/strain_panel.py
\begin{table}[htbp]
\centering
\caption{\textbf{Three of the paper's headline strains are one construction step past a tube already in the freezer.} Each row is a single transformation, and in two cases a FACS enrichment, away from a strain we hold.}
\label{tab:gap}
\footnotesize
\begin{tabular}{@{}ll >{\raggedright\arraybackslash}p{62mm} l@{}}
\toprule
Strain & From & Step & Titer (mg/L) \\
\midrule
YZy312 (Strain B) & C0052 YZy91 & transform with pYZ228, the CEN URA3 plasmid carrying PTDH3-ILV6V110E-TADH1; this completes the Fig. 3c 2x2 that C0043, C0044 and C0045 are three quarters of & -- \\
YZy505 & C0053 YZy502 & transform with the 2$\mu$ plasmid pYZ350 (two copies of Ec\_ilvCP2D1-A1, one each of Ll\_IlvDI433V, \gene{ARO10}, Ll\_adhARE1), then two rounds of FACS & 681 $\pm$ 29 \\
YZy470 & C0051 YZy469 & move Ll\_ilvDI433V from the CEN plasmid to a 2$\mu$ plasmid & 443 $\pm$ 6 \\
YZy148 + \gene{LEU4} mutant \#6 & C0046 YZy148 & transform with the evolved \gene{LEU4} plasmid (E86D K191N K374R A445T S481R N515I A568V S601A) & 963 $\pm$ 32 \\
\bottomrule
\end{tabular}
\end{table}


\notesec{The green group is three quarters of a 2$\times$2} YZy311, YZy313 and
YZy314 (C0043--C0045) are named nowhere in the paper's running text, which is
why the tube labels give nothing away. Supplementary Table 1 places all three at
once, and Fig.~3c is where they appear:

\begin{quote}\small
Isobutanol titers of YZy91 transformed with CEN plasmids containing ILV6$^{WT}$
(Strain A) or ILV6$^{V110E}$ (\gene{ILV6} mutant \#6, Strain B) compared to
titers from a strain overexpressing the mitochondrial isobutanol pathway
(YZy363) transformed with CEN plasmids containing ILV6$^{WT}$ (Strain C) or
ILV6$^{V110E}$ (\gene{ILV6} mutant \#6, Strain D).
\end{quote}

\noindent
That is a 2$\times$2: two hosts crossed with two \gene{ILV6} alleles. C0043 is
Strain A, C0044 is Strain C, and C0045 is Strain D. The fourth cell, Strain B,
is YZy312, and it is not in the box.

%% SOURCE: experiments/034-isobutanol-wetlab/scripts/strain_panel.py
%% Mirror: $DATA_ROOT/torchcell-library/zhangBiosensorBranchedchainAmino2022/si/si3.pdf
%% sha256 74de9b6cff654c5d9b6988a0a0c08672eb7defad4f002ecce6b9e6b9fa48824f
\begin{table}[htbp]
\centering
\caption{\textbf{Genotypes, as Supplementary Table 1 prints them.} Parent names the strain each was built from, so a chain can be walked back to CEN.PK2-1C (MATa \gene{ura3}-52 \gene{trp1}-289 \gene{leu2}-3,112 \gene{his3}-1 MAL2-8c SUC2). Plasmid contents are in parentheses, following the source.}
\label{tab:genotypes}
\footnotesize
\begin{tabular}{@{}lll >{\raggedright\arraybackslash}p{88mm}@{}}
\toprule
ID & Strain & Parent & Genotype \\
\midrule
C0043 & Yzy311 & YZy91 & YZy91, CEN URA3 plasmid (PTDH3-ILV6-TADH1) \\
C0044 & YZy313 & YZy363 & YZy363, CEN URA3 plasmid (PTDH3-ILV6-TADH1) \\
C0045 & YZy314 & YZy363 & YZy363, CEN URA3 plasmid (PTDH3-ILV6V110E-TADH1) \\
C0046 & YZy148 & CEN.PK2-1C & CEN.PK2-1C, his3::HIS3-PLEU1-yEGFP-TADH1 \gene{bat1}$\Delta$::hphMX \gene{leu4}$\Delta$::lox71-kanMX-lox66 \gene{leu9}$\Delta$::lox71-natMX-lox66 leu2::LEU2 \\
C0047 & YZy148+LEU4WT & YZy148 & YZy148, CEN plasmid (PTDH3-LEU4-TADH1), wild-type allele \\
C0048 & YZy148+LEU4deltaS547 & YZy148 & YZy148, CEN plasmid carrying \gene{LEU4}$\Delta$S547 \\
C0049 & YZy452 & YZy449 & YZy449, $\delta$-integration-PTDH3-Ec\_ilvCP2D1-A1-TADH1-[PTEF1-Ec\_ilvCP2D1-A1-TACT1]. YZy449 is \gene{ilv3}$\Delta$ \gene{tma29}$\Delta$ with PGAL10-Ll\_ilvD and PTEF1-Bs\_alsS. \\
C0050 & YZy454 & YZy452 & YZy452, CEN URA3 plasmid (PTDH3-Ll\_ilvD-TADH1) \\
C0051 & YZy469 & YZy452 & YZy452, CEN URA3 plasmid (PTDH3-Ll\_ilvDI433V-TADH1) \\
C0052 & YZy91 & CEN.PK2-1C & CEN.PK2-1C, his3::HIS3-PLEU1-yEGFP-PEST-TADH1-PTPI1-LEU41-410-TPGK1, \gene{bat1}$\Delta$::hphMX \gene{bat2}$\Delta$::lox71-kanMX-lox66 \gene{ilv6}$\Delta$::lox71-natMX-lox66 \\
C0053 & YZy502 & YZy487 & YZy487, $\delta$-integration-PC120-PDC1-TACT1-PTDH3-Ec\_ilvCP2D1-A1-TCYC1-PTEF1-Ll\_ilvDI433V-TTPS1-PGAL1-S-Bs\_alsS-TACT1 \\
\bottomrule
\end{tabular}
\end{table}


\notesec{All eleven are yeast, and seven of them carry a plasmid} Nothing in
this panel is an \org{E.\ coli} strain. Every one of the eleven appears in
Supplementary Table 1, which the source titles "Yeast strains used in this
study", and the only \org{E.\ coli} in the paper is DH5$\alpha$ as a cloning
host, plus \gene{Ec\_ilvC}, an \org{E.\ coli} gene expressed in yeast.

The distinction that does cut through the panel is episomal. C0046, C0049, C0052
and C0053 are chromosomal only. The other seven each carry one plasmid, and all
seven plasmids are the same backbone, AmpR CEN \gene{URA3}, built off the empty
parent pYZ125. Three things follow, and none of them is visible in a genotype
string.

\begin{itemize}\setlength\itemsep{0.15em}
  \item \textbf{They have to be held on uracil dropout.} A CEN \gene{URA3}
    plasmid under no selection segregates, and a plasmid-bearing tube streaked
    on rich medium quietly becomes its own parent. The paper's own methods grow
    these in "SC (or SC-ura)" for exactly this reason.
  \item \textbf{\gene{URA3} is spent in all seven.} That is the marker a guide
    library wants, so none of the seven can host one without curing the plasmid
    first.
  \item \textbf{A plasmid is a separate request from a strain.} AmpR means the
    plasmids live in \org{E.\ coli} in the source lab. Asking for pYZ228 is not
    the same transaction as asking for a yeast strain, and it is the cheaper of
    the two when the combination can be rebuilt here.
\end{itemize}

\noindent
The last point simplifies two items on the request list. pYZ125 is the empty
vector, so the missing vector control is C0049 transformed with pYZ125 rather
than a request for YZy453; and YZy312 is C0052 transformed with pYZ228 rather
than a request for a strain.

\notesec{Titers down that column do not compare} Four measurement conditions
appear in \cref{tab:panel}: 15 percent galactose, 15 percent glucose, 2 percent
glucose in the dark, and a 48 hour fermentation scored for a different product.
What survives that is the set of contrasts in which exactly one thing changes,
and there are three of them.

\begin{itemize}\setlength\itemsep{0.15em}
  \item \textbf{C0050 against C0051}, isobutanol, 15 percent glucose.
    \gene{Ll\_ilvD} wild type against \gene{Ll\_ilvD}$^{\text{I433V}}$ on the
    same CEN backbone in the same host. 3.1-fold, the largest single-allele
    effect in the panel.
  \item \textbf{C0044 against C0045}, isobutanol. \gene{ILV6} wild type against
    \gene{ILV6}$^{\text{V110E}}$, same host, 1.6-fold.
  \item \textbf{C0047 against C0048}, isopentanol. \gene{LEU4} wild type
    against \gene{LEU4}$\Delta$S547, same host, with C0046 as the allele-free
    floor beneath both.
\end{itemize}

\noindent
C0043 does not pair with C0044, although both carry wild-type \gene{ILV6}:
their hosts differ by the \gene{bat2}$\Delta$ \gene{ilv6}$\Delta$ deletions and
by the mitochondrial pathway at once, so the two differ in more than one thing.

\notesec{There is no wild type in the box} Every one of the eleven carries the
biosensor cassette, and most carry deletions in \gene{BAT1}, \gene{BAT2},
\gene{ILV6}, \gene{LEU4}, \gene{LEU9} or the three \gene{PDC} genes. An
unengineered CEN.PK2-1C reference is not among them, and a differential
expression analysis wants one. It is a cheap thing to request alongside the
next shipment.

%% notes-tex/034-isobutanol-wetlab/sections/2-yko.tex
%% SOURCE: hand-authored. Every quote is verbatim from two mirrored documents,
%%   $DATA_ROOT/torchcell-library/zhangBiosensorBranchedchainAmino2022/  (the paper)
%%   $DATA_ROOT/torchcell-library/lopezSystemsMetabolicEngineering2024/  (the dissertation)
%% No number in this section is derived.

\section{The Knockout-Collection Work Shares a Cassette, Not a Background\secstatus{todo}}
\label{sec:yko}

These strains get grouped with the isobutanol screen of the yeast gene-knockout
collection (\term{YKOC}) that ran in the same lab, and the grouping is half
right. The screen is chapter 3 of Jos\'e de Jes\'us Monta\~no L\'opez's
dissertation, \textit{Systems Metabolic Engineering of Isobutanol Production in
Saccharomyces cerevisiae} (Princeton, November 2024, adviser Jose L.\ Avalos),
mirrored as \file{lopezSystemsMetabolicEngineering2024}. It names the paper
above as its reference 158:

\begin{quote}\small
We used a previously designed branched-chain amino acid biosensor to rapidly
quantify flux changes in the isobutanol biosynthetic pathway$^{158}$
\end{quote}

\notesec{The cassette is the same construct, to the promoter} The dissertation's
Supplementary Table 5 gives the sensor strain as

\begin{quote}\small
YKOC, \gene{his3}::HIS3-P$_{\text{LEU1}}$-yEGFP-PEST-T$_{\text{ADH1}}$-P$_{\text{TPI1}}$-LEU4$_{1-410}$-T$_{\text{PGK1}}$
\end{quote}

\noindent
which is the paper's isobutanol configuration term for term: destabilized yEGFP
read out of P$_{\text{LEU1}}$, amplified by the truncated LEU4$_{1-410}$ under
P$_{\text{TPI1}}$. The \gene{leu2}$\Delta$ the configuration calls for comes free,
because every strain in the knockout collection is \gene{leu2}$\Delta$0. So a
strain carrying the isobutanol configuration in our box and a strain from that
screen are reporting the same quantity through the same reporter.

\notesec{The background is where they part} The screen runs in BY4741, which the
same table gives as \term{S288C} MATa \gene{his3}$\Delta$1 \gene{leu2}$\Delta$0
\gene{met15}$\Delta$0 \gene{ura3}$\Delta$0. The paper runs in CEN.PK2-1C. These
are different laboratory lineages, and the paper's own choice of CEN.PK2-1C for
higher alcohol work is itself the argument that the choice of lineage is not
neutral for this phenotype.

The consequence for sequencing is specific. A differential expression contrast
between one of our tubes and a strain from the knockout screen would confound
the perturbation with the lineage, and no normalization recovers that. The two
sets support a comparison of \emph{direction}, not of expression level.

\notesec{None of the eleven is a strain from that screen} The dissertation names
no strain with a \texttt{YZy} prefix anywhere in its text, and the paper names
none of the dissertation's \texttt{yJM} strains. The nearest relative to the
screen in our box is any tube carrying the isobutanol configuration, and what it
shares is the reporter.

\notesec{What to ask for, if the screen background is the point} The strain the
screen was anchored on is \term{yJM1}, BY4741 carrying that same cassette at
\gene{HIS3}, and the collection itself is the biosensor-containing YKOC. Both
are Avalos-lab strains and neither is in the shipment. If the transcriptomics
plan is meant to connect to the knockout screen rather than to the biosensor
paper, yJM1 is the request to make, and it is a different request from the one
that produced these eleven tubes.

%% notes-tex/034-isobutanol-wetlab/sections/3-rnaseq.tex
%% SOURCE: hand-authored, resting on tables/t1-strain-panel.tex,
%% tables/t3-genotypes.tex and the two mirrored documents cited in
%% sections/1-panel.tex and sections/2-yko.tex.
%% The design fact that drives the selection, that Ll_ilvD in YZy452 sits behind
%% P_GAL10 and is therefore off in glucose, is quoted below.

\section{The First Sequencing Round\secstatus{todo}}
\label{sec:rnaseq}

The ask is a low-flux anchor, the best producer, and steps in between. The panel
supplies two independent versions of that, both for isobutanol, and one sentence
in the Methods is what turns the first of them into three points rather than
two.

\notesec{The promoter that makes a ladder} The dehydratase in YZy452 is not
constitutive.

\begin{quote}\small
Because the wild-type copy of \gene{Ll\_ilvD} in YZy452 is expressed from a
P$_{\text{GAL10}}$ promoter, growing the library in glucose ensures this
wild-type copy is repressed, and only the \gene{Ll\_ilvD} variants from the CEN
plasmid library (inserted downstream of a P$_{\text{TDH3}}$ promoter) are
expressed.
\end{quote}

\noindent
YZy452 also carries \gene{ilv3}$\Delta$, so the native dihydroxyacid dehydratase
is gone. In glucose, C0049 therefore has no dehydratase at all, C0050 supplies
the wild-type bacterial enzyme from a CEN plasmid, and C0051 supplies the
improved allele from the same backbone. Three levels of one step, one host, one
carbon source.

\notesec{Round one, five strains} Two matched contrasts, on two different steps
of the pathway: \gene{Ll\_ilvD} wild type against I433V at 3.1-fold, and
\gene{ILV6} wild type against V110E at 1.6-fold, with C0049 as the floor beneath
the first. \Cref{tab:roundone} gives the five, and what every tube needs if it is
picked up at all.

%% SOURCE: experiments/034-isobutanol-wetlab/scripts/strain_panel.py
%% Growth conditions from the Methods of zhangBiosensorBranchedchainAmino2022:
%% single colony -> 1 mL SC or SC-ura + 2% glucose overnight -> 10 uL into
%% 1 mL of the same in a 24-well plate -> 20 h -> spin, resuspend in 1 mL
%% of the same + 15% sugar -> seal -> 48 h, 30 C, 200 rpm -> OD600, HPLC.
\begin{table}[htbp]
\centering
\caption{\textbf{Round one, and how each tube is grown.} Medium is uracil dropout exactly when a CEN \gene{URA3} plasmid has to be held, and the fermentation carbon source is not a free choice for C0049. Every strain marked yes runs on the paper's ordinary 24-well protocol: single colony, shaker, plate reader, HPLC, with no sorting and no special illumination. Leucine must be present in all cases, because \gene{leu2}$\Delta$ makes these strains leucine auxotrophs; standard SC supplies it.}
\label{tab:roundone}
\footnotesize
\begin{tabular}{@{}lllll >{\raggedright\arraybackslash}p{50mm}@{}}
\toprule
ID & Strain & Medium & Fermentation & Pick and use & Round one \\
\midrule
C0044 & YZy313 & SC-ura & 15\% glucose & yes & yes -- \gene{ILV6} wild type, low arm of the 1.6-fold pair \\
C0045 & YZy314 & SC-ura & 15\% glucose & yes & yes -- ILV6V110E, high arm of the 1.6-fold pair \\
C0049 & YZy452 & SC & 15\% galactose & yes & yes -- ladder floor, dehydratase off on glucose \\
C0050 & YZy454 & SC-ura & 15\% glucose & yes & yes -- wild-type allele anchor \\
C0051 & YZy469 & SC-ura & 15\% glucose & yes & yes -- Ll\_ilvDI433V, 3.1-fold over C0050 \\
\midrule
C0043 & Yzy311 & SC-ura & 15\% glucose & yes & -- \\
C0046 & YZy148 & SC & 15\% glucose & yes & -- \\
C0047 & YZy148+LEU4WT & SC-ura & 15\% glucose & yes & -- \\
C0048 & YZy148+LEU4deltaS547 & SC-ura & 15\% glucose & yes & -- \\
C0052 & YZy91 & SC & 15\% glucose & yes & -- \\
C0053 & YZy502 & SC & 2\% glucose, dark & \textbf{no} & -- \\
\bottomrule
\end{tabular}
\end{table}


\notesec{No sorting anywhere in this} The protocol that produced every titer in
the paper is a single colony into 1 mL of SC or SC-ura with 2 percent glucose
overnight, 10 $\mu$L of that into 1 mL of the same in a 24-well plate, 20 h, then
a spin and a resuspension into 1 mL of the same medium with 15 percent sugar,
sealed, 48 h at 30 $^{\circ}$C and 200 rpm, then OD$_{600}$ on a plate reader and
HPLC on the supernatant. A shaker, a plate reader and an HPLC. Fluorescence-activated
cell sorting appears only in the path that isolates strains out of pooled
libraries, and those isolations are already done.

\notesec{Two conditions that are not free choices} C0049 ferments on galactose,
not glucose, for the P$_{\text{GAL10}}$ reason above, and on glucose it will read
as a dead strain rather than as a floor. C0053 is the one tube that is not pick
and use at all: its overnight needs constant blue light before production can run
in the dark, which is an illuminated-plate setup rather than a shaker.

Leucine has to be in the medium for all of them. The SI is direct about why:
"LEU2 deletion also causes leucine auxotrophy, requiring supplementation of this
amino acid in the growth medium". Standard SC supplies it, and a leucine dropout
would kill the panel.

\notesec{Measure titer alongside, do not carry it across} C0049's published
310 $\pm$ 15 mg/L is a galactose number. Whichever carbon source the round
settles on, the titer for each strain belongs in the same experiment as the
sequencing rather than taken from the paper, because the media here will not
match the media there well enough for the numbers to transfer.

\noindent
C0053 stays out of round one on the light grounds above, but it is worth
sequencing eventually for one reason: it is the direct parent of YZy505, the
paper's best isobutanol strain at 681 $\pm$ 29 mg/L, so a baseline taken now is
what makes that comparison possible later.

\notesec{What is not in the round} The isopentanol arm, C0046 through C0048, is
the cleanest three-point series in the box and C0048 holds the highest titer of
anything here at 842 $\pm$ 23 mg/L. It is the leucine branch. This program is
the valine branch, so those three are round two, and they move up only if a
dose-response in $\alpha$-isopropylmalate flux is wanted for its own sake.

C0043 YZy311 sits out because its nearest neighbor in the panel, C0044, differs
from it by both host and pathway, so a contrast between them would not isolate
anything. C0052 YZy91 sits out because it is the plasmid-free host of C0043
rather than a producer.

\notesec{Three requests worth making with the next shipment}

\begin{itemize}\setlength\itemsep{0.15em}
  \item \textbf{YZy121}, which Supplementary Table 1 gives as CEN.PK2-1C
    carrying the isobutanol-configured biosensor and nothing else. It is the
    reference every strain in this panel is missing: same lineage, same
    reporter, no deletions and no pathway.
  \item \textbf{YZy312}, Strain B, which completes the Fig.~3c 2$\times$2 that
    C0043, C0044 and C0045 are three quarters of. It is one transformation of
    C0052 with pYZ228.
  \item \textbf{YZy453}, YZy452 carrying an empty CEN URA3 plasmid. C0049
    carries no plasmid while C0050 and C0051 do, so the plasmid itself and its
    selection marker are currently confounded with the dehydratase.
\end{itemize}

\noindent
None of the three connects this panel to the knockout screen; for that reason
\cref{sec:yko} names yJM1 instead.


\clearpage
\small
\bibliography{references}

\end{document}
